\section{Detailed Analysis}

\subsection{Validity Analysis of Execution Outcomes}
\label{app:validation}

\subsubsection{Motivation of Validity Analysis}

In the realm of artificial intelligence and machine learning, the efficacy, precision, and reliability of models are crucial for practical implementations. Evaluating multiple models provides an understanding of their respective strengths and limitations, leading to better informed decisions about which models are best suited for specific tasks. The purpose of this validity analysis is to offer a systematic approach to discern how different models perform, particularly in terms of task completion, context size constraints, return format accuracy, action accuracy, and task limitations. This deep dive into performance parameters not only enhances our knowledge about the models' capabilities, but also aids in refining and optimizing them for future applications.

\subsubsection{Definition of Validity Analysis}

For comprehensive validity analysis, we have demarcated the results into five distinct categories:

\begin{itemize}
    \item \textbf{Completed}: Denotes instances where models, irrespective of the end outcome, successfully finished the task as per the instructions.
    
    \item \textbf{Context Limit Exceeded}: Denotes instances where the model's length was constrained by the API, predominantly observed in the \texttt{text-davinci} model.
    
    \item \textbf{Invalid Format}: Denotes instances where models, despite receiving clear instructions, failed to return responses in the expected format.
    
    \item \textbf{Invalid Action}: Denotes instances where the models returned in the correct format, but their actions either fell outside the permitted action space or had incorrect action parameters.
    
    \item \textbf{Task Limit Exceeded}: Denotes instances tasks reached their termination criteria, such as exceeding the stipulated number of rounds.
\end{itemize}

By categorizing the results into these classes, we can gain a clearer picture of where each model excels and where they encounter challenges, allowing for targeted improvements.

\subsubsection{Validity Analysis of Models}

For our evaluation, we scrutinized the validity performance of \num distinct models. Apart from the \texttt{text-davinci} model, which has an inherent strict API context length constraint, the outcomes for other models primarily fall under the categories of Completed, Invalid Format, Invalid Action, and Task Limit Exceeded.

From the detailed analysis showcased, key trends emerge. As depicted in Figure~\ref{fig:agent_validation}, the chart offers a clear visualization of the validity distribution across distinct models and defined categories, enabling us to derive insightful conclusions.


\begin{figure}[t]
    \centering
    \caption{Validity analysis of models. \textbf{Invalid Format}, \textbf{Invalid Action}, and \textbf{Text Limit Exceeded} are common errors. \textbf{Context Limit Exceeded} errors only appear in \texttt{text-davinci} models.}
    \includegraphics[width=\linewidth]{figs/agent_validation.pdf}
    \label{fig:agent_validation}
    \vspace{-5mm}
\end{figure}